\section{Ход работы}

\subsection{Схема цепи (пункт~\ref{task:scheme})}
Схема исследуемой цепи показана на рисунке~\ref{fig:scheme}. Генератор синусоидального
напряжения $u(t)$ питает последовательно соединённые резистор $R$ и конденсатор $C$.
Амперметр включён последовательно и измеряет ток цепи, вольтметр — параллельно
конденсатору. Оба прибора показывают действующие значения.

\bbtScreenshot[0.75\linewidth]{materials/schema.png}{Схема последовательной
RC-цепи}{fig:scheme}

\subsection{Предварительный расчёт (пункт~\ref{task:calc})}
Сопротивление конденсатора зависит от частоты по формуле~\eqref{eq:xc}:
\begin{equation}\label{eq:xc}
  X_C = \frac{1}{2\pi f C},
\end{equation}
где $f$ — частота, Гц;\\
$C$ — ёмкость, Ф.

Полное сопротивление, ток и напряжение на конденсаторе — по формулам~\eqref{eq:z}
и~\eqref{eq:i}:
\begin{equation}\label{eq:z}
  Z = \sqrt{R^2 + X_C^2},
\end{equation}
\begin{equation}\label{eq:i}
  I = \frac{U}{Z}, \qquad U_C = I X_C,
\end{equation}
где $R$ — сопротивление резистора, Ом;\\
$U$ — действующее напряжение генератора, В.

Ток опережает напряжение генератора на угол $|\varphi|$, где
\[
  \varphi = -\operatorname{arctg} \frac{X_C}{R}.
\]

Например, на частоте 1000~Гц:
\begin{gather*}
  X_C = \frac{1}{2\pi \cdot 1000 \cdot 0{,}22 \cdot 10^{-6}} \approx 723{,}4~\text{Ом}, \\
  Z = \sqrt{1500^2 + 723{,}4^2} \approx 1665{,}3~\text{Ом}, \qquad
  I = \frac{5}{1665{,}3} \approx 3{,}00~\text{мА}.
\end{gather*}
Результаты расчёта для всех частот приведены в таблице~\ref{tab:calc}.

\begin{bbttable}{Результаты предварительного расчёта}{tab:calc}
  \begin{tabular}{|r|r|r|c|c|c|}
    \hline
    $f$, Гц & $X_C$, Ом & $Z$, Ом & $I$, мА & $U_C$, В & $\varphi$, ° \\ \hline
    100  & 7234,3 & 7388,2 & 0,68 & 4,90 & −78,3 \\ \hline
    200  & 3617,2 & 3915,8 & 1,28 & 4,62 & −67,5 \\ \hline
    500  & 1446,9 & 2084,1 & 2,40 & 3,47 & −44,0 \\ \hline
    1000 & 723,4  & 1665,3 & 3,00 & 2,17 & −25,7 \\ \hline
    2000 & 361,7  & 1543,0 & 3,24 & 1,17 & −13,6 \\ \hline
    5000 & 144,7  & 1507,0 & 3,32 & 0,48 & −5,5  \\ \hline
  \end{tabular}
\end{bbttable}

\subsection{Измерения (пункт~\ref{task:measure})}
Ток и напряжение на конденсаторе измерены на стенде на шести частотах из задания при
напряжении генератора 5~В. Показания приборов приведены вместе с расчётом в
таблице~\ref{tab:compare}.

\subsection{Сравнение расчёта и эксперимента (пункт~\ref{task:compare})}
В таблице~\ref{tab:compare} расхождение $\delta$ — разность эксперимента и расчёта в
процентах от расчёта.

\begin{bbttable}{Расчёт и эксперимент}{tab:compare}
  \begin{tabular}{|r|c|c|c|c|c|c|}
    \hline
    & \multicolumn{3}{c|}{Ток $I$, мА} & \multicolumn{3}{c|}{Напряжение $U_C$, В} \\
    \cline{2-7}
    $f$, Гц & расчёт & эксп. & $\delta$, \% & расчёт & эксп. & $\delta$, \% \\ \hline
    100  & 0,68 & 0,71 & 4,9 & 4,90 & 4,89 & −0,1 \\ \hline
    200  & 1,28 & 1,33 & 4,2 & 4,62 & 4,59 & −0,6 \\ \hline
    500  & 2,40 & 2,47 & 3,0 & 3,47 & 3,40 & −2,1 \\ \hline
    1000 & 3,00 & 3,05 & 1,6 & 2,17 & 2,10 & −3,3 \\ \hline
    2000 & 3,24 & 3,27 & 0,9 & 1,17 & 1,13 & −3,6 \\ \hline
    5000 & 3,32 & 3,35 & 1,0 & 0,48 & 0,46 & −4,2 \\ \hline
  \end{tabular}
\end{bbttable}

Зависимости тока и напряжения на конденсаторе от частоты показаны на
рисунке~\ref{fig:graph}: линии — расчёт, точки — эксперимент.

\bbtScreenshot[0.85\linewidth]{materials/grafik.png}{Зависимости $I(f)$ и $U_C(f)$, расчёт и
эксперимент}{fig:graph}

\subsection{Анализ результатов (пункт~\ref{task:explain})}
С ростом частоты сопротивление конденсатора по формуле~\eqref{eq:xc} падает, поэтому ток
растёт и стремится к $U/R \approx 3{,}33$~мА, а напряжение на конденсаторе падает. Сдвиг фаз
уменьшается по модулю с −78,3 до −5,5°: на низких частотах цепь ведёт себя почти как
ёмкость, на высоких — почти как резистор. Граница между этими режимами — частота, на которой
$X_C = R$:
\[
  f_\text{с} = \frac{1}{2\pi R C} \approx 482~\text{Гц}.
\]
Около неё, на частоте 500~Гц, напряжения на резисторе и конденсаторе близки, а
$U_C \approx U / \sqrt{2}$.

Эксперимент повторяет расчётные зависимости, но расхождение меняется с частотой закономерно,
а не случайно:
\begin{itemize}
  \item ток на низких частотах выше расчётного на 4,2–4,9~\%, а на высоких — всего на
  0,9–1,0~\%;
  \item напряжение на конденсаторе, наоборот, почти совпадает с расчётом на низких частотах
  (−0,1~\%) и ниже расчётного на высоких (до −4,2~\%).
\end{itemize}
На низких частотах ток определяется в основном сопротивлением конденсатора, на высоких —
резистора. Значит, сопротивление конденсатора на стенде меньше расчётного, то есть его
ёмкость больше номинальной примерно на 5~\%. Это в пределах допуска ±10~\%. На высоких
частотах $U_C \approx U X_C / R$, поэтому та же ошибка ёмкости уменьшает и напряжение на
конденсаторе. Расхождение тока на высоких частотах около 1~\% соответствует допуску
резистора ±1~\% и точности отсчёта приборов до сотых.
